\begin{table}[htbp]
\centering
\caption{Intended-use specification for the two prediction models. Both are \emph{research risk-stratification} tools only; neither is a screening, diagnostic, or triage device. The specification is fixed before validation so that performance is judged against a declared purpose.}
\label{tab:intended}
\small
\begin{tabular}{p{0.24\linewidth}p{0.70\linewidth}}
\toprule
Dimension & Specification \\
\midrule
Purpose & Research risk stratification for cohort description and hypothesis generation; \emph{not} individual screening, diagnosis, triage, or treatment guidance. \\
Target user & Researchers/epidemiologists analyzing survey or, later, deidentified cohort data. \\
Prediction time & Cross-sectional (concurrent survey visit); no temporal/prognostic claim. \\
Target population & Adults 40--69 y in the fixed audiometry windows; not generalized beyond. \\
Inputs & Prespecified harmonized covariates (Table~\ref{tab:variables}); no free-text or image inputs in the public-data models. \\
Output / decision & Calibrated probability for research stratification; drives \emph{no} clinical action. \\
Cost of false positive & Low in research use; would be non-trivial if misused for screening (over-referral) --- such use is contraindicated. \\
Cost of false negative & Low in research use; dangerous if misused clinically (missed disease) --- such use is contraindicated. \\
Contraindicated uses & Individual diagnosis, treatment decisions, insurance/employment decisions, or any deployment without prospective clinical validation and regulatory review. \\
Safety interaction & The deterministic red-flag layer (Table~\ref{tab:redflag}) overrides model output for time-critical symptoms regardless of predicted probability. \\
\bottomrule
\end{tabular}
\end{table}
